\begin{annexitem}{annex:p12.har.fw-normalization}{thm:p12-har-fw-normalization}{The weak LF host and normal-head inversion}
\begin{proof}
The fixed signature translates into the PTS context by reading the LF
kind type as \(\ast\), its ambient kind classifier as \(\Box\),
and preserving object/family products, abstractions and applications.
Every object-domain product has one of the two specified LF triples.
The given signatures are well formed in order: prop,A,D use the sort
axiom; object constructors/relations use the star/star product;
T and eq use the star/Box family product; every proof constant uses
star/star products over its declared indices and evidence, ending in
the corresponding T/eq family. The inherited T declaration is retained
in PRED. No kind-domain product is needed for any listed declaration.
Both triples and the axiom are in exact CC; constants are fixed context
entries with the same types. The seven typing-rule induction cases of
\cref{thm:p12-har-pts-morphism} therefore give CC typing for every
LF judgment. \Cref{thm:p12-har-cc-normalization} supplies strong beta
normalization, including annotations, without a closed-context premise.

Typed reduction remains in LF, rather than merely in its larger CC
target. Here is the rule construction. In the combined PTS view fixed
signature entries may have star or Box certificates; actual LF object
variable extensions have a star certificate, never a kind annotation.
Projection follows by the
variable rule at that prefix and weakening through the suffix. Thinning
into a valid context containing the literal old declarations is proved
by induction: sort axioms weaken through the new certificates; a variable
uses projection; a weakening is absorbed by the containing context;
products/abstractions first transport their annotation, then their body
into the valid fresh-binder extension; application combines its two
induction judgments; conversion combines its subject and typed result.
These exhaust the seven rules, and also justify the premise transports
used in the following generation argument. Generation strips terminal
weakening/conversion to obtain variable, product, abstraction and
application witnesses. A conversion composes their type equality; a
weakening transports their premises into the valid enlarged context.
Product injectivity follows from the raw common-reduct argument because
beta cannot change a product's outer constructor. Type correctness in
application obtains the product's codomain sort premise by generation,
then substitutes the actual typed argument using
\cref{thm:p12-har-pts-substitution}. These operations use only the
two LF triples and the axiom.

Context conversion is derived by induction: at the changed declaration,
type the variable at its new typed annotation and convert to the old
typed one; transform each later annotation before reusing variable or
weakening. Product/abstraction/application/conversion rebuild their
original rule and its same admitted LF pair. For beta preservation,
induct on the subject. At a root application of a lambda, product
injectivity converts the argument to the literal binder annotation,
and substitution gives the body's instance type. Type correctness
supplies the application's result-type sort for the restoring conversion.
For reductions inside a binder annotation, use its smaller-subject
induction, convert the binder context, and rebuild the same product or
lambda. Body and application-child reductions use their smaller-subject
inductions. If the original type is the untyped top \(\Box\),
no final conversion to it is permitted: no judgment types that atom.
Variable/lambda/application subjects at that literal result are excluded
by their formation/type-correctness premises; the remaining sort or
kind-product cases retain their literal result sort. This covers the
top exception without an illicit conversion judgment.

Thus a finite normalizing path gives a beta-normal LF subject of its
original type. In a beta-normal term, a base-typed object is a variable
or constant followed by an application spine: a root lambda has product
type, not an opaque atomic family type, and a lambda-headed application
would be a redex. Generation determines every argument type in that
spine. Atomic families remain headed by their opaque names; products
remain products under every beta reduct. Raw confluence consequently
prevents conversion between them. This gives the claimed normal-head
inversion, not a beta-conversion of every term to an eta-long form.
\end{proof}
\end{annexitem}

\begin{annexitem}{annex:p12.har.fw-reification}{thm:p12-har-fw-reification}{All normal-name grammars, binding cases, and open substitutions}
\begin{proof}
Normalize object and formula names with the preceding lemma; normalization
is unique up to bound-name changes by raw beta confluence. Inspect normal
head types in the exact signatures. In PROP, a prop-typed head is one
of the declared atom names or imp, with exactly its two prop arguments.
Proof variables have types headed by T, not prop, and proof constants
return T-families; none can create a prop name. Hence the stated formula
grammar follows inductively on its two proper arguments.

In PRED, a normal A-typed head is an object variable, c, or f with one
A argument. A normal prop-typed head is R with two such terms, imp with
two prop names, or forall with its A-to-prop function. Normalization and
head-type inversion rule out every T/proof head in these categories.
A normal function of type A-to-prop is a lambda with prop-typed body,
or a neutral head with a spine. Inspection of the signature leaves
only R with one A argument in the latter case: f's result is A;
imp takes prop rather than A; forall takes an A-to-prop function rather
than A; contexts have no other function variables. The opaque A and
prop families are distinct normal forms, so conversion cannot change
these arity/domain comparisons.

Consequently reifying a forall-lambda recurses on its smaller body,
after replacing its binder by a fresh variable of the same type.
Reifying forall(Rs) directly gives the stated atomic body under a fresh
quantifier; it does not call itself again on a term of the same size.
Imp recurses on its smaller children, R on its data arguments, and f
on its smaller data child. These cases prove definedness and termination,
including the eta-short family case.

For the lambda signature a normal D-typed head is a free D variable,
app with two D arguments, or abs with its D-to-D function. No equality
proof or family can supply a D-typed head. A normal D-to-D function
is a lambda with D body, or app applied to one D argument. The context
contains no function variables, and abs expects a function rather than
a D argument. Thus the abs-lambda branch recurses on its smaller body,
and the eta-short abs(app u) branch recurses on the proper subterm u,
then constructs the source lambda/application. This exhausts its grammar.

Freshen every binder against the finite context and substituted images.
Induct on these grammars to prove substitution commutation. A replaced
object variable gives its decoded image; another variable or constant
is unchanged. The first-order f and R cases and the lambda app case
combine their child inductions. In a forall- or abs-lambda case the
fresh binder permits substitution in its body, and reification constructs
the corresponding capture-avoiding source binder. In an eta-short
function case its original object argument is substituted and the new
fresh point variable is not; the child induction therefore gives the
same decoded application under the binder. Substitution of a base A
or D name cannot introduce a host lambda in function position: the
replaced variables have atomic type, not product type. The resulting
normal-name grammar is preserved; any preceding beta reductions commute
with substitution by the general term substitution identity. This gives
the stated equalities for arbitrary open object substitutions.

Induction on source expressions proves the retraction on their encodings:
variables/constants and applications use the corresponding clause, and
an encoded binder is abs/forall applied to its explicitly abstracted
encoded body, so fresh-variable beta application recovers that body.
Finally any LF beta-convertible names have identical normal forms;
normalization followed by this reifier therefore gives the same decoded
source expression up to binder renaming. This is conversion invariance,
not an object eta rule or a claimed bijection on every LF term.
\end{proof}
\end{annexitem}

\begin{annexitem}{annex:p12.har.fw-faithfulness}{thm:p12-har-fw-faithfulness}{Actual inverse derivations for PROP, PRED, and beta equality}
\begin{proof}
Normalize the given proof term, retaining its type by
\cref{thm:p12-har-fw-normalization}. Decode each assumption index
by \cref{thm:p12-har-fw-reification}. For PROP and PRED prove the
slightly more general assertion: a normal proof of T(q) constructs a
minimal-logic derivation of the decoded normal name q from the decoded
assumption names in the current context. Newly introduced proof-variable
types may name an intermediate formula rather than a literal original
source encoding; their decoded assumptions are carried explicitly.

The proof construction is well founded on the size of the normal
evidence term, counted with its annotations. At an introduction requiring
a function-valued proof f, use a fresh variable z of its stated domain
and normalize fz. If f is a lambda, this is its smaller body with
its bound variable renamed to z; substitution of a fresh variable creates
no new redex. If f is neutral, fz has size at most size(f)+2.
The enclosing imp-intr term has f plus a head, three applications and
two formula parameters; forall-intr has f plus a head, two applications
and its family parameter. In either case the new evidence is strictly
smaller than the enclosing evidence. All elimination evidence arguments
are already proper subterms. This justifies the recursive construction
also for eta-short proof functions; it is not an assumed canonicalization.

In PROP, normal proof-head inversion leaves exactly three possibilities.
A proof variable h_j gives the corresponding decoded assumption.
At imp-intr p q f, introduce fresh h:T(p), decode the normalized fh
to a derivation of the decoded q under the added decoded p assumption,
then discharge h by implication introduction. At imp-el p q u v,
recursively decode u as an implication from decoded p to decoded q
and v as decoded p; implication elimination gives decoded q.
Atomic head/type injectivity identifies these decoded conclusions with
the required target. No other signature head returns a T-family.

PRED has these three cases and two additional ones. At forall-intr F f,
introduce fresh z:A and decode the normalized fz. Its proof type is
T(Fz), so its conclusion is the reified pointwise body. The fresh z
occurs in none of the old object declarations or assumption indices;
the reifier adds no new free variables. Universal introduction therefore
has the actual eigenvariable condition and concludes the decoded
forall F. At forall-elim F u t, decode u to that universal formula,
and decode the actual normal object argument t to its first-order term.
Universal elimination yields the reified body with this decoded term
substituted. Reification/substitution commutation identifies it with
the decoded normal form of Ft, exactly the evidence's result index.
Thus the proof works with open arguments, not only with ground c-terms.
Extra proof variables introduced by implication remain encoded assumptions;
they cannot create object/formula names by the preceding grammar argument.

For the lambda signature, strengthen the construction to every normal
proof of eq u v in a free-D-variable context: it produces a derivation
of beta equality between the decoded u and v. Normal proof-head inversion
has no proof-variable case here, because the admissible context has only
D variables. The complete constant cases are as follows.

At refl u use reflexivity. At sym u v e decode e and use symmetry.
At trans u v w e d decode both evidence arguments and use transitivity.
At mon x x' z z' e d, decode e for the argument equality and d for
the function equality, then apply source application congruence to
obtain decoded z applied to decoded x equal to decoded z' applied
to decoded x'. This is the source mon argument order, not a swap of
its two proof inputs.

At xi F G e, introduce fresh z:D and normalize ez. The same size
argument applies: the enclosing xi includes both F and G parameters,
so ez is strictly smaller even if e is neutral. Decode it to
\(\mathcal D(\mathrm{nf}(Fz))=_\beta
\mathcal D(\mathrm{nf}(Gz))\), then use source abstraction congruence
at that fresh variable. Its two abstractions are precisely the decoded
abs F and abs G. No equation assumption depending on z was introduced.

At beta F u, the left endpoint decodes to
\((\lambda z.\mathcal D(\mathrm{nf}(Fz)))\mathcal D(u)\).
One source beta step contracts it to the decoded body substituted by
\(\mathcal D(u)/z\). The reification/substitution identity identifies
this with \(\mathcal D(\mathrm{nf}(Fu))\), the right endpoint.
For a lambda F this is ordinary body substitution; for F=app v it
is the application \(\mathcal D(v)\mathcal D(u)\). These are
both possible normal function forms, so no object eta law was used.
All other source equational steps are generated by these cases.

For every signature, a generated result type may beta-convert to the
given result type. Normalize its indices: opacity of T or eq and raw
confluence make the normal indices agree. The reifier therefore gives
the same decoded conclusion. At an original encoded target the
retraction lemma gives the original formula or lambda term. Combining
with the already supplied forward rule inductions gives both existence
directions under precisely the stated signature/context conditions.
\end{proof}
\end{annexitem}

\begin{annexitem}{annex:p12.har.fw-no-eta}{thm:p12-har-fw-no-eta}{A concrete nonderivable eta equation}
\begin{proof}
If the stated LF inhabitant existed, the lambda converse would give
\(u=_\beta\lambda x.ux\) in the untyped calculus. Both sides
are beta-normal and not alpha-equivalent. For completeness, raw
confluence applies to this untyped claim without a typing premise:
annotate every source lambda binder by a fixed sort atom s. This
injective translation preserves every beta contraction and beta-normal
form; its image is closed under reduction since s does not reduce.
The two translated terms would therefore be convertible normal forms
in the raw grammar of \cref{thm:p12-har-pseudo-confluence}, contradicting
its common-reduct property. Hence no inhabitant exists. This distinguishes
pointwise reification from adding object eta equality.
\end{proof}
\end{annexitem}

\begin{annexitem}{annex:p12.har.rational-normalization}{thm:p12-har-rational-normalization}{Gcd reduction, Bezout cancellation, and the witness-sensitive carrier}
\begin{proof}
Write n=p+1 and m=q+1, so every denominator is positive. Cross-product
equality is reflexive and symmetric. If am=bn and bk=cm, multiply
the first equation by k and the second by n, then cancel the positive
m to obtain ak=cn. Thus it is an equivalence relation.

The positive common divisors of a and n form a nonempty finite subset
of \(\{1,\ldots,n\}\); define d as their maximum. This definition
handles negative a by its absolute value and gives d=n when a=0.
Both a/d and n/d are integers, with n/d positive. If k>1 divided
both reduced coordinates, kd would be a positive common divisor of
a,n larger than d, a contradiction. Thus the reduced pair is coprime.
Its cross-product with the original pair is identical after multiplying
by d. In particular zero reduces to numerator0, denominator1, namely
the source pair(0,0).

Prove the cancellation fact required for uniqueness rather than assuming
it. For coprime a,n, the integer combinations ua+vn include a positive
number n; let e be their least positive member. Integer division by e
exists directly by finite search: for a nonnegative integer choose the
largest q with qe at most that integer; its remainder is below e.
For a negative integer, negate this division, decreasing its quotient
by one when the positive remainder is nonzero. Divide a and n in
this way. Each remainder is an integer combination and is between0
and e-1; a positive one would contradict minimality. Hence e divides
both a and n. Coprimality forces e=1, so there are integers u,v with
ua+vn=1. If n divides am, multiply that identity by m; both terms
on the right are divisible by n, and consequently n divides m.

Now suppose reduced pairs a/n,b/m have am=bn. Applying this fact
to a,n gives n divides m, and applying it to b,m gives m divides n.
Their positivity implies n=m; cancellation then gives a=b. This
proves uniqueness of the numerical reduced pair in each equivalence
class, including zero and negative numerators.

To construct r into the source subsetoid, pair each computed reduced
integer/natural pair with the coprimality proof just given. The original
and reduced pairs are equal in the source relation. If two originals
are equivalent, their reduced numerical pairs are identical by uniqueness,
so r preserves the subsetoid equality. Projection j preserves equality
by its definition, and jr is pointwise equal to the original by the
cross-product calculation. At a coprime pair d=1, so the first component
of rj is that pair. Its possibly newly constructed coprimality proof
need not be the original proof, but the subsetoid equality ignores both
witnesses. Hence rj is pointwise equal to its identity as well.
Nothing here proves that the dependent-pair proof witnesses themselves
are definitionally equal or unique inhabitants.
\end{proof}
\end{annexitem}

\begin{annexitem}{annex:p12.har.setoid-quotient-universal}{thm:p12-har-setoid-quotient-universal}{Both map directions and uniqueness in the declared equality}
\begin{proof}
The identity q preserves the old equality because it is included in Q.
If a function f from A to B preserves Q, pair it with precisely that
proof to obtain its map from [A,Q]. Its composite with q has the same
underlying function as f, hence is equal pointwise in the target setoid.
It is also a map from the old setoid, by precomposing the preservation
proof with the inclusion of the old equality in Q.

Conversely any map from [A,Q] has, by its defining type, a preservation
proof for Q, so its composite with q is an old-setoid map with that
extra property. These constructions are inverse up to function-setoid
equality: they retain the same value at every argument and disregard
the possibly different preservation proof fields. If another descended
map has the same composite with q, that equality at each a:A compares
exactly the two descended values, since q is the identity on the carrier.
It is therefore pointwise equal to the first descended map. This proves
the full mapping property without a choice of quotient representatives.
\end{proof}
\end{annexitem}
